% articulo-06-romanos.tex - Numeración romana de páginas preliminares.
%
% Tres variantes controladas desde el frontmatter YAML:
%   romanos-minusculas: true  → i, ii, iii… (norma anglosajona, Chicago 18.ª § 1.4)
%   romanos-mayusculas: true  → I, II, III… (RAE, Ortografía § 2.1.1h)
%   (sin declarar)            → ɪ, ɪɪ, ɪɪɪ… versalitas (compromiso tipográfico, por defecto)
%
% \@roman es el comando interno que LaTeX llama para producir \roman{N}.
% Redefinirlo aquí afecta solo a los contextos de numeración romana
% (páginas preliminares); la numeración arábiga del cuerpo no se ve afectada.

\makeatletter
$if(romanos-minusculas)$
% Minúsculas — sin redefinición; \pagenumbering{roman} produce minúsculas por defecto.
$else$$if(romanos-mayusculas)$
\renewcommand{\@roman}[1]{\MakeUppercase{\romannumeral#1}}
$else$
% Versalitas (por defecto)
\renewcommand{\@roman}[1]{\textsc{\romannumeral#1}}
$endif$$endif$
\makeatother
